\labsection{5}{Interprocess communication and synchronization}{sec:lab05}
\begin{labproject}{Compare a pipe, shared memory, and producer-consumer state}

\textbf{Coverage.} Chapter~\ref{ch:ipc}: IPC models, shared state, synchronization conditions.\\
\textbf{Goal.} Move one message through a pipe and another through a shared-memory segment.\\
\textbf{Explain.} Explain who owns each communication object and why shared memory needs separate synchronization.

\section*{Source code}
\labsourcecard{\texttt{pipe\_demo.c}}{Parent writes bytes; child reads them through a unidirectional pipe.}{Each process closes the pipe end it does not use. Closing the write end is also how end-of-stream can be signalled.}
\lstinputlisting[style=utpcnumbered]{all_experiments/Lab05/pipe_demo.c}

\labsourcecard{\texttt{producer\_consumer\_sim.c}}{Make full/empty buffer state visible without threads yet.}{The invariant is \texttt{empty + full = CAPACITY}; blocked operations do not change the counters.}
\lstinputlisting[style=utpcnumbered]{all_experiments/Lab05/producer_consumer_sim.c}

\labsourcecard{\texttt{shm\_writer.c}}{Create/attach a System V shared-memory segment and place a message in it.}{The writer detaches but deliberately leaves the segment available for another process.}
\lstinputlisting[style=utpcnumbered]{all_experiments/Lab05/shm_writer.c}

\labsourcecard{\texttt{shm\_reader.c}}{Attach to the existing segment, read the message, then mark the segment for removal.}{The reader demonstrates cleanup with \texttt{shmdt()} and \texttt{shmctl(..., IPC\_RMID,...)}.}
\lstinputlisting[style=utpcnumbered]{all_experiments/Lab05/shm_reader.c}

\section*{Task 1 --- Pipe: follow the data path}
\begin{lstlisting}[style=utpbash]
cd all_experiments/Lab05
gcc -Wall -Wextra -std=c11 pipe_demo.c -o pipe_demo
./pipe_demo
\end{lstlisting}
Draw one arrow from parent to kernel pipe to child. Label which file descriptor end each process closes.

\section*{Task 2 --- Producer/consumer state}
\begin{lstlisting}[style=utpbash]
gcc -Wall -Wextra -std=c11 producer_consumer_sim.c -o producer_consumer_sim
./producer_consumer_sim
\end{lstlisting}
For every row, verify \texttt{empty + full = 3}. Circle the first blocked produce and first blocked consume. Explain why the counters do not change on a blocked action.

\section*{Task 3 --- Shared-memory lifecycle}
\begin{lstlisting}[style=utpbash]
gcc -Wall -Wextra -std=c11 shm_writer.c -o shm_writer
gcc -Wall -Wextra -std=c11 shm_reader.c -o shm_reader
./shm_writer
ipcs -m
./shm_reader
ipcs -m
\end{lstlisting}
Record the segment ID after the writer and check whether it remains after the reader marks it for removal.

\section*{Task 4 --- Failure case}
Run \texttt{./shm\_reader} again after the segment has been removed. The failure is intentional. Explain why a clear error message is better than creating a new empty segment in the reader.

\begin{quickcheck}
\begin{enumerate}
  \item How is a pipe different from shared memory?
  \item What makes a section of code a critical section?
  \item Why can two correct operations still produce an incorrect result when they race?
\end{enumerate}
\end{quickcheck}

\begin{verification}
\begin{itemize}
  \item Each process closes the pipe end it does not use, and the sketch shows which.
  \item Every producer/consumer row satisfies \texttt{empty + full = 3}, and blocked actions leave the counters unchanged.
  \item The shared-memory segment ID is recorded after the writer, and its removal is confirmed with \texttt{ipcs -m}.
  \item The pipe and shared memory are compared on data movement, lifetime and synchronization.
\end{itemize}
\end{verification}

\begin{labdeliverable}
Submit: pipe sketch, producer-consumer state check, shared-memory before/after evidence, and a comparison table with columns \textbf{pipe} vs. \textbf{shared memory}: data movement, lifetime, and synchronization concern.
\end{labdeliverable}
\end{labproject}
